% GENERATED FILE. Edit canonical lessons, then run npm run generate:books.
% canonical-source-hash: fnv1a64:3bbe507033c344ce
% canonical-lessons: RU-C359-obyazannost, RU-C359-poruchat, RU-C359-poruchenie, RU-C359-vypolnyat, RU-C359-spravlyatsya

\chapter{Duties at work}
\label{ch:ru-duties-at-work}
\hlchaptermodality{359}

\begin{chapteropening}
\textbf{By the end of this chapter:} \emph{I can talk about my duties at work: what I am given to do, what I carry out and whether I can cope.}

Complete the last lesson of chapter 359: I can talk about my duties at work: what I am given to do, what I carry out and whether I can cope.
\end{chapteropening}

\section[\texorpdfstring{\ru{обязанность} (obyázannost)}{обязанность (obyázannost)}]{\ru{обязанность} (obyázannost) — a duty, a responsibility}

\label{lesson:RU-C359-obyazannost}



\begin{quote}
Before the new one: say the Russian for to hire, then the Russian for length of service.
\end{quote}

\subsection*{You'll want to know: \ru{обязанность}}
\textbf{\ru{обязанность}} (\emph{obyázannost}) — \textquotedblleft{}a duty, a responsibility\textquotedblright{}.

\begin{quote}
\ru{В} \ru{мои} \ru{обязанности} \ru{входит} \ru{работа} \ru{с} \ru{клиентами}.
\end{quote}

\begin{quote}
\emph{V moí obyázannosti vkhódit rabóta s kliyéntami.}
\end{quote}

\begin{quote}
My duties include working with clients.
\end{quote}

Feminine. \emph{\ru{Входить} \ru{в} \ru{обязанности}} — to be part of one’s duties.

\begin{cousinweb}
\textbf{\ru{обязанность}} is built on \emph{\ru{обязать}} \textquotedblleft{}to oblige\textquotedblright{}, made of \emph{\ru{о}-} and \emph{\ru{вязать}} \textquotedblleft{}to tie, to bind\textquotedblright{}, from Proto-Slavic \emph{*vęzati}: a duty is what binds you.
\end{cousinweb}

\begin{grammarlens}[title={What you've built}]
One more word for telling, explaining and describing at B1.
\end{grammarlens}

\subsection*{Your turn}
\begin{itemize}
\raggedright
  \item \emph{Say it:} \emph{obyázannost}
  \item \emph{Say it:} \emph{obyázannost}, once more
  \item \emph{Say it:} a sentence of your own with \emph{obyázannost}
  \item \emph{Recall:} say the Russian for to hire, then the Russian for length of service, then say \emph{obyázannost} again
\end{itemize}

\begin{tcolorbox}[breakable,colback=teal!4,colframe=teal!35!black,title={Before you move on}]
What does \ru{обязанность} mean? (\textquotedblleft{}A duty, a responsibility\textquotedblright{}.) Say it once more.
\end{tcolorbox}
\section[\texorpdfstring{\ru{поручать} (poruchát)}{поручать (poruchát)}]{\ru{поручать} (poruchát) — to entrust, to put in charge of}

\label{lesson:RU-C359-poruchat}



\begin{quote}
Before the new one: say the Russian for length of service, then the Russian for a duty.
\end{quote}

\subsection*{You'll want to know: \ru{поручать}}
\textbf{\ru{поручать}} (\emph{poruchát}) — \textquotedblleft{}to entrust, to put in charge of\textquotedblright{}.

\begin{quote}
\ru{Ему} \ru{часто} \ru{поручают} \ru{самые} \ru{трудные} \ru{задачи}.
\end{quote}

\begin{quote}
\emph{Yemú chásto porucháyut sámyye trúdnyye zadáchi.}
\end{quote}

\begin{quote}
He is often given the hardest tasks.
\end{quote}

Imperfective; perfective \emph{\ru{поручить}} (\emph{\ru{поручу}, \ru{поручишь}}). \emph{\ru{Поручать} \ru{кому} \ru{что}}, or + infinitive.

\begin{cousinweb}
\textbf{\ru{поручать}} \textquotedblleft{}to commission, to entrust\textquotedblright{} is built on \emph{\ru{ручать}} \textquotedblleft{}to vouch, to guarantee\textquotedblright{}, from \emph{\ru{рука}} \textquotedblleft{}hand\textquotedblright{} (EtymDB). To entrust a task is to put it into someone’s hands.
\end{cousinweb}

\begin{grammarlens}[title={What you've built}]
One more word for telling, explaining and describing at B1.
\end{grammarlens}

\subsection*{Your turn}
\begin{itemize}
\raggedright
  \item \emph{Say it:} \emph{poruchát}
  \item \emph{Say it:} \emph{poruchát}, once more
  \item \emph{Say it:} a sentence of your own with \emph{poruchát}
  \item \emph{Recall:} say the Russian for length of service, then the Russian for a duty, then say \emph{poruchát} again
\end{itemize}

\begin{tcolorbox}[breakable,colback=teal!4,colframe=teal!35!black,title={Before you move on}]
What does \ru{поручать} mean? (\textquotedblleft{}To entrust, to put in charge of\textquotedblright{}.) Say it once more.
\end{tcolorbox}
\section[\texorpdfstring{\ru{поручение} (poruchéniye)}{поручение (poruchéniye)}]{\ru{поручение} (poruchéniye) — an assignment, an errand}

\label{lesson:RU-C359-poruchenie}



\begin{quote}
Before the new one: say the Russian for a duty, then the Russian for to entrust.
\end{quote}

\subsection*{You'll want to know: \ru{поручение}}
\textbf{\ru{поручение}} (\emph{poruchéniye}) — \textquotedblleft{}an assignment, an errand\textquotedblright{}.

\begin{quote}
\ru{Начальник} \ru{дал} \ru{мне} \ru{важное} \ru{поручение}.
\end{quote}

\begin{quote}
\emph{Nachálnik dal mne vázhnoye poruchéniye.}
\end{quote}

\begin{quote}
The boss gave me an important assignment.
\end{quote}

Neuter. \emph{\ru{Выполнить} \ru{поручение}} — to carry out an assignment.

\begin{cousinweb}
\textbf{\ru{поручение}} is derived from \emph{\ru{поручать}} \textquotedblleft{}to entrust\textquotedblright{}, from \emph{\ru{ручать}} \textquotedblleft{}to vouch\textquotedblright{}, from \emph{\ru{рука}} \textquotedblleft{}hand\textquotedblright{} (EtymDB): a task placed in someone’s hands.
\end{cousinweb}

\begin{grammarlens}[title={What you've built}]
One more word for telling, explaining and describing at B1.
\end{grammarlens}

\subsection*{Your turn}
\begin{itemize}
\raggedright
  \item \emph{Say it:} \emph{poruchéniye}
  \item \emph{Say it:} \emph{poruchéniye}, once more
  \item \emph{Say it:} a sentence of your own with \emph{poruchéniye}
  \item \emph{Recall:} say the Russian for a duty, then the Russian for to entrust, then say \emph{poruchéniye} again
\end{itemize}

\begin{tcolorbox}[breakable,colback=teal!4,colframe=teal!35!black,title={Before you move on}]
What does \ru{поручение} mean? (\textquotedblleft{}An assignment, an errand\textquotedblright{}.) Say it once more.
\end{tcolorbox}
\section[\texorpdfstring{\ru{выполнять} (vypolnyát)}{выполнять (vypolnyát)}]{\ru{выполнять} (vypolnyát) — to carry out, to fulfil}

\label{lesson:RU-C359-vypolnyat}



\begin{quote}
Before the new one: say the Russian for to entrust, then the Russian for an assignment.
\end{quote}

\subsection*{You'll want to know: \ru{выполнять}}
\textbf{\ru{выполнять}} (\emph{vypolnyát}) — \textquotedblleft{}to carry out, to fulfil\textquotedblright{}.

\begin{quote}
\ru{Он} \ru{всегда} \ru{выполняет} \ru{свои} \ru{обещания}.
\end{quote}

\begin{quote}
\emph{On vsegdá vypolnyáyet svoí obeshchániya.}
\end{quote}

\begin{quote}
He always keeps his promises.
\end{quote}

Imperfective; perfective \emph{\ru{выполнить}}. \emph{\ru{Выполнять} \ru{работу}, \ru{план}, \ru{просьбу}}.

\begin{cousinweb}
\textbf{\ru{выполнять}} is built from \emph{\ru{вы}-} \textquotedblleft{}out\textquotedblright{} and the root of \emph{\ru{полный}} \textquotedblleft{}full\textquotedblright{}, inherited from Proto-Slavic \emph{*pĭlnŭ} (EtymDB): to fill a task out to the end.
\end{cousinweb}

\begin{grammarlens}[title={What you've built}]
One more word for telling, explaining and describing at B1.
\end{grammarlens}

\subsection*{Your turn}
\begin{itemize}
\raggedright
  \item \emph{Say it:} \emph{vypolnyát}
  \item \emph{Say it:} \emph{vypolnyát}, once more
  \item \emph{Say it:} a sentence of your own with \emph{vypolnyát}
  \item \emph{Recall:} say the Russian for to entrust, then the Russian for an assignment, then say \emph{vypolnyát} again
\end{itemize}

\begin{tcolorbox}[breakable,colback=teal!4,colframe=teal!35!black,title={Before you move on}]
What does \ru{выполнять} mean? (\textquotedblleft{}To carry out, to fulfil\textquotedblright{}.) Say it once more.
\end{tcolorbox}
\section[\texorpdfstring{\ru{справляться} (spravlyátsya)}{справляться (spravlyátsya)}]{\ru{справляться} (spravlyátsya) — to cope, to manage}

\label{lesson:RU-C359-spravlyatsya}



\begin{quote}
Before the new one: say the Russian for an assignment, then the Russian for to carry out.
\end{quote}

\subsection*{You'll want to know: \ru{справляться}}
\textbf{\ru{справляться}} (\emph{spravlyátsya}) — \textquotedblleft{}to cope, to manage\textquotedblright{}.

\begin{quote}
\ru{Я} \ru{не} \ru{справляюсь} \ru{с} \ru{работой} \ru{один}.
\end{quote}

\begin{quote}
\emph{Ya ne spravlyáyus s rabótoy odín.}
\end{quote}

\begin{quote}
I cannot cope with the work on my own.
\end{quote}

Imperfective; perfective \emph{\ru{справиться}} (\emph{\ru{справлюсь}, \ru{справишься}}). It takes \emph{\ru{с}} + instrumental.

\begin{cousinweb}
\textbf{\ru{справляться}} is built on \emph{\ru{править}} \textquotedblleft{}to set straight, to steer\textquotedblright{}, from \emph{\ru{правый}} \textquotedblleft{}right\textquotedblright{}, inherited from Proto-Slavic \emph{*pravŭ} \textquotedblleft{}right, correct\textquotedblright{} (EtymDB). To cope is to set a thing straight.
\end{cousinweb}

\begin{grammarlens}[title={What you've built}]
One more word for telling, explaining and describing at B1.
\end{grammarlens}

\subsection*{Your turn}
\begin{itemize}
\raggedright
  \item \emph{Say it:} \emph{spravlyátsya}
  \item \emph{Say it:} \emph{spravlyátsya}, once more
  \item \emph{Say it:} a sentence of your own with \emph{spravlyátsya}
  \item \emph{Recall:} say the Russian for an assignment, then the Russian for to carry out, then say \emph{spravlyátsya} again
\end{itemize}

\begin{tcolorbox}[breakable,colback=teal!4,colframe=teal!35!black,title={Before you move on}]
What does \ru{справляться} mean? (\textquotedblleft{}To cope, to manage\textquotedblright{}.) Say it once more.
\end{tcolorbox}
